\documentclass[output=paper,colorlinks,citecolor=brown,
% hidelinks,
% showindex
]{langscibook}

\newcommand{\citeapos}[1]{\citeauthor{#1}'s \citeyearpar{#1}}
\newcommand{\laszAsp}[0]{[asp]}
\newcommand{\laszLoc}[1]{[loc:#1]}
\newcommand{\laszDir}[2]{[dir:#1$\to$#2]}
\newcommand{\laszHs}[1]{[hs:#1]}
\newcommand{\laszClaw}[0]{{\fontspec{aslfontgithubio.ttf} S}}
\newcommand{\laszPlain}[0]{[plain]}

\author{
    Sabine Laszakovits\affiliation{University of Connecticut; Austrian Academy of Sciences} and
    Ronice Müller de Quadros\affiliation{Universidade Federal de Santa Catarina} and 
    {Emily Jo} Noschese\affiliation{University of Hawai'i} and 
    Diane Lillo-Martin\affiliation{University of Connecticut; Haskins Laboratories}
}
\title{Object shift in ASL and Libras}
\abstract{
    ASL and Libras have an object-shift construction 
    by which the canonical SVO order is changed to SOV. 
    In both sign languages, this ordering is 
    mandatory for V marked with durative/continuative aspect (reduplicated movement), 
    optional for V that agrees with O in locus,
    and not allowed with plain V. 
    When V agrees with O in handshape, 
    ASL requires OV ordering whereas Libras allows both OV and VO ordering. 
    We present an analysis that derives these data 
    with a combination of syntactic movement of O 
    and violable, equally-ranked PF-constraints as proposed by \citet{BW.2012}. 
    Unlike \citeapos{Matsuoka.1997} and \citeapos{Braze.2004} proposals, 
    we do not move V to a head on the right in violation of the Final-over-final constraint
    \citep{BiberauerHR.2014,SheehanBRH.2017}.
}

\begin{document}
\maketitle

% ===================================================================
\section{Introduction}

The underlying word order in modern American Sign Language (ASL)
and Brazilian sign language (Língua Brasileira de Sinais, Libras) has
often been argued to be subject-verb-object (SVO; 
\citealp[see]{Fischer.1975,Liddell.1980,Padden.1988} for ASL; 
\citealp{Quadros.1999,Quadros.2003} for Libras), 
which we adopt here. However, in both languages, word order
variations are possible. In this paper, we focus on the construction with
the word order SOV, which has been termed ``object shift''. This
construction differs from topicalization structures like O'SV, VO'S,
S'VO, or S'O'V in that the latter contain a prosodic break after the
topic (indicated here by an apostrophe) and may contain nonmanual
marking (raised brows and a slight upward chin tilt; 
\citealp{Liddell.1977,Liddell.1980,Padden.1988}) during the utterance of the topic.%
\footnote{
    \citet{Fischer.1990} argues that object-shift constructions in ASL involve
    ``mini-topicalization'' of O to some intermediate specifier position. She corroborates
    this analysis by reporting a definiteness effect found in SOV orders, which is typical
    of topics. However, unlike typical topicalizations, SOV constructions lack a prosodic
    break after O and cannot mark O with nonmanual topicalization marking. Our
    analysis is similar to \citeauthor{Fischer.1990}'s in that we think O moves to some intermediate
    specifier position. However, we have not been able to reproduce this definiteness
    effect for ASL, and we see no reason to call this movement topicalization.
}
We restrict our attention to verbs that take a direct object (such as `buy X') and
possibly an indirect object (such as `send X to Y'), but do not discuss
verbs that denote movement to or from a location (such as `put X on Z').%
\footnote{
    Furthermore, we distinguish object-shift from locative constructions of the
    form location-subject-predicate, which have been taken as instances of the word
    order OSV because, like object-shift, they lack the prosodic break of topicalization
    structures \citep{Liddell.1980}. Locative constructions differ from object-shift
    constructions at least in the following properties: 
    \begin{inparaenum}[(i)]
        \item word order is OSV (vs. SOV),
        \item the predicate is restricted to `is located at' (vs. any verb), and 
        \item there is no minimally different derivation with the word order SVO (which there is for
    object-shift; see e.g.~\REF{lasz:ex:4}--\REF{lasz:ex:7} and throughout the paper).
    \end{inparaenum}
}

There are three main triggers for object-shift in ASL and
Libras: 
\begin{inparaenum}[1.]
    \item durative/continuative aspect on the verb, which makes object-shift obligatory; 
    \item verbal agreement in handshape with the object, which makes object-shift 
    optional in Libras and obligatory in ASL; and 
    \item verbal agreement in locus with the object, which makes object-shift optional in both languages. 
\end{inparaenum}
The effect of durative aspect on word order is illustrated in \REF{lasz:ex:1}--\REF{lasz:ex:2}.%
\footnote{
    Following standard practice in sign linguistics, signs are glossed using
    English words in all caps. Most of our examples represent sign sequences that are
    grammatical in both ASL and Libras, though the actual signs are of course different.
    `IX' is the gloss used for a pointing indexical sign which serves
    pronominal functions. `\laszAsp{}' indicates an aspectual marker involving reduplication
    of the sign root. `\laszHs{\_}' indicates the use of a classifier handshape. `\laszLoc{\_}'
    indicates that a sign is produced using a spatial locus other than default (`a', `b',
    etc. indicate distinct loci but not physical location). `\laszDir{\_}{\_}' indicates that a sign
    moves from one locus to another. \par 
    In the English translation, past tense is used although neither ASL nor
    Libras has grammatical tense marking; similarly, we often use `a/the' for nouns
    because in both of these sign languages either translation is possible.
}%
\footnote{
    When a source is not provided, the judgments primarily come from the
    2\textsuperscript{nd} and 3\textsuperscript{rd} co-authors, who are native signers of Libras and ASL respectively.
}

\ea
    \label{lasz:ex:1}
    \ea[]{
        \label{lasz:ex:1a}
        IX1 WINE DRINK\laszAsp{}. 
        \hfill 
        \cmark ASL, \cmark Libras \\ 
        `I drank wine continuously.'
    }
    \ex[*]{
        \label{lasz:ex:1b}
        IX1 DRINK\laszAsp{} WINE. 
        \hfill 
        *ASL, *Libras
    }
    \z 
\ex
    \label{lasz:ex:2}
    \ea[]{
        \label{lasz:ex:2a}
        MY SISTER LETTER SEND\laszAsp{}. 
        \hfill 
        \cmark ASL, \cmark Libras \\ 
        `My sister repeatedly sent a/the letter(s).'
    }
    \ex[*]{ 
        \label{lasz:ex:2b}
        MY SISTER SEND\laszAsp{} LETTER. 
        \hfill 
        *ASL, *Libras 
    }
    \z 
\z

Verbal agreement can target the object's noun-class%
\footnote{
    In this sense we follow \citeapos{Benedicto.Brentari.2004} proposal, whereby
    classifiers are treated as a type of agreement between V and O/S.
}, 
whereby V
changes its handshape (`hs') to a classifier as illustrated in~\REF{lasz:ex:3}, or it
can target the object's referent, whereby V changes the end-point of its
movement (`dir') to the object's locus (`loc'), as seen in~\REF{lasz:ex:4}.%
\footnote{
    \citet{Fischer.1975} proposed that object-shift is also licensed if the relation
    between V and O is established in the semantics, such that SOV is possible only if it
    is clear from the semantics of V, S, and O which argument is the agent and which is
    the patient. However, \citet{Liddell.1980} showed that this approach makes incorrect
    predictions, and that syntactic agreement is a better way to go.
}

\ea 
    \label{lasz:ex:3}
    The handshape of `GIVE' changes to the \laszClaw{} handshape indicative of handling an apple. 
    \ea[]{ 
        SALLY APPLE GIVE\laszHs{\laszClaw}. 
        \hfill 
        \cmark ASL, \cmark Libras 
    }
    \ex[]{ 
        SALLY GIVE\laszHs{\laszClaw} APPLE. 
        \hfill
        \cmark Libras 
    }
    \ex[*]{
        SALLY GIVE\laszHs{\laszClaw} APPLE. 
        \hfill 
        *ASL \\ 
        all: `Sally gave someone a/the apple.'
    }
    \z 
\ex
    \label{lasz:ex:4}
    The direction of `HELP' changes to the locus of `ANA'. 
    \ea 
        IX\laszLoc{a} MARIA\laszLoc{a} IX\laszLoc{b} ANA\laszLoc{b} HELP\laszDir{a}{b}. \\
        \citep[from][6]{Quadros.etal.2004}
        \hfill 
        \cmark Libras, \cmark ASL 
    \ex 
        IX\laszLoc{a} MARIA\laszLoc{a} HELP\laszDir{a}{b} IX\laszLoc{b} ANA\laszLoc{b}. \\
        \citep[from][5]{Quadros.etal.2004} 
        \hfill 
        \cmark Libras, \cmark ASL \\
        both: `Maria helped Ana.'
    \z 
\z 
 
For some verbs, locus agreement involves signing the verb at the
object's locus as illustrated in \REF{lasz:ex:5} (termed ``locationality'' by \citealp{Fischer.Gough.1978}, 
``spatialization'' by \citealp{Quadros.etal.2004}, 
and ``co-localization'' by \citealp{Lourenco.Wilbur.2018}; 
see also \citealp{Bergman.1980,Liddell.1980,Costello.2015,Smith.1990}, i.a.; ASL \& Libras),
or by adding an auxiliary as in \REF{lasz:ex:6} (Libras only; \citealp{Quadros.1999}). We
will treat all of these options as locus-agreement strategies and
postulate a uniform syntax.

\ea 
    \label{lasz:ex:5}
    Spatialized plain verbs
    \ea 
        MAN BICYCLE\laszLoc{a} BUY\laszLoc{a}. 
        \hfill 
        \cmark ASL, \cmark Libras \\ 
        \citep[from][9]{Quadros.etal.2004}
    \ex 
        MAN BUY\laszLoc{a} BICYCLE\laszLoc{a}. 
        \hfill 
        \cmark ASL, \cmark Libras \\ 
        \citep[from][9]{Quadros.etal.2004} \\ 
        both: `The man bought a/the bicycle.'
    \z 
\ex 
    \label{lasz:ex:6}
    Auxiliary 
    \ea 
        IX\laszLoc{a} JOÃO\laszLoc{a} IX\laszLoc{b} MARIA\laszLoc{b} AUX\laszDir{a}{b} SUPPORT. 
        \hfill 
        \cmark Libras 
        \\ 
        \citep[from][7]{Quadros.etal.2004} 
    \ex 
        IX\laszLoc{a} JOÃO\laszLoc{a} SUPPORT IX\laszLoc{b} MARIA\laszLoc{b}. 
        \hfill 
        \cmark Libras 
        \\ 
        \citep[from][5]{Quadros.etal.2004}
    \z 
\z 

The object-shift construction is not possible if neither trigger is
present, i.e., when the verb is ``plain'', as in \REF{lasz:ex:7}.

\ea 
    \label{lasz:ex:7}
    \ea[*]{
        MAN NUMBER FORGET\laszPlain{}. 
        \hfill 
        *ASL
        \\ 
        \citep[from][]{Liddell.1980}
    }
    \ex[*]{
        IX JOHN SOCCER LIKE\laszPlain{}. 
        \hfill 
        *Libras
        \\ 
        \citep[from][61]{Quadros.1999}
    }
    \ex[]{
        MAN FORGET\laszPlain{} NUMBER. 
        \hfill 
        \cmark ASL 
        \\ 
        `The man forgot a/the number.'
        \citep[from][]{Liddell.1980}
    }
    \ex[]{
        IX JOHN LIKE\laszPlain{} SOCCER. 
        \hfill 
        \cmark Libras 
        \\ 
        `John likes soccer.'
        \citep[from][61]{Quadros.1999}
    }
    \z 
\z 

Table~\ref{lasz:tab:1} presents a summary of the acceptability judgments for the
basic order SVO and for the object-shift construction SOV for both
ASL and Libras.

\newcommand{\laszTabColA}[0]{\cellcolor{gray!30}}
\newcommand{\laszTabColB}[0]{\cellcolor{gray!10}}
\newcommand{\laszTabColC}[0]{\cellcolor{gray!30}}
% \newcolumntype{C}[1]{>{\centering\let\newline\\\arraybackslash\hspace{0pt}}m{#1}} % aligns in the middle vertically

\begin{table}
    \centering 
    \begin{tabular}{r c c c c}
        \hline
        \hline 
        & V\laszPlain{} & V\laszLoc{\_} & V\laszHs{\_} & V\laszAsp{} \\ 
        \hline 
        ASL & 
            \laszTabColA{} & 
            \laszTabColB{} & 
            \laszTabColC{} & 
            \laszTabColC{} 
            \\ 
        Libras & 
            \multirow{-2}{*}{\laszTabColA{} \shortstack{\cmark SVO \\ \xmark SOV}} & 
            \multirow{-2}{*}{\laszTabColB{} \shortstack{\cmark SVO \\ \cmark SOV}} & 
            \laszTabColB{} & 
            \multirow{-2}{*}{\laszTabColC{} \shortstack{\xmark SVO \\ \cmark SOV}} 
            \\
        \hline 
    \end{tabular}
    \caption{Comparison of judgments for SVO and OSV}
    \label{lasz:tab:1}
\end{table}

Several proposals have been made regarding the syntax of object-shift
in ASL/Libras. The difficulty in providing a satisfactory analysis of
this phenomenon seems to stem from the disparate nature and the
disparate obligatoriness of the constructions that trigger object shift.

For durative-aspect marking of V, \citet{Matsuoka.1997} and \citet{Braze.2004} 
have proposed that object shift arises because aspectual
morphology sits on an Asp head, and V must move there, and AspP is
right-headed. In section~\ref{lasz:sec:3} we will show that there is an empirical and a
conceptual problem with this analysis.

An alternative proposal is based on the fact that durative-aspect
marking is expressed by reduplication of the verb’s movement 
\citep{Klima.Bellugi.1979}. The complexity of a sign’s movement corresponds to
its phonological heaviness \citep{Brentari.1998}, and we know from several
spoken languages that there is a preference for phonologically heavy
elements to appear late in a sentence \parentext{\citealp{Ross.1967,Williams.2003}, i.a.}.
\citet{Brentari.1998} has argued that this preference also applies in the
object-shift construction with heavy verbs, but remains agnostic
whether the inversion is realized by V moving to the right, or by O
moving to the left triggered by V’s heaviness, as \citet{Liddell.1980}
proposed.%
\footnote{
    Note that according to \citeapos{Brentari.1998} analysis, locus-agreement per se
    does not constitute `complex movement', neither if it modifies the locus in which V
    is signed, nor if it modifies the locus of V’s endpoint or starting point. In section 6.5
    she discusses how phonological weight units interact with the order of syntactic
    constituents. In \citeauthor{Brentari.1998}’s model, movement corresponds to prosodic features, and
    each branching prosodic feature constitutes a weight unit. In ASL, GIVE usually
    contains only a “direction”-movement, that is one weight unit. Both VO and OV
    ordering are allowed. When habitual aspect is added, GIVE contains an additional
    “trill” movement, a total of two weight units. OV ordering is preferred over VO.
    Exhaustive aspect adds both a “trill” movement and an “arc” movement, totallying
    three weight units. Now, due to V’s heaviness, VO is unacceptable.
}

However, as we have seen, object shift is also related to
agreement between V and O, both noun-class (classifier) agreement
\REF{lasz:ex:3}; and referential (locus) agreement \REF{lasz:ex:4}--\REF{lasz:ex:5}. 
These agreement-data
are difficult to unify with the aspect-data. If one follows a syntactic
proposal for the aspect-data like \citet{Matsuoka.1997} or \citet{Braze.2004},
one has to explain why aspect-marking makes the OV order
obligatory, while agreement-marking allows both the OV and the VO
order. On the other hand, if one follows a phonology-driven proposal
for the aspect-data like \citet{Liddell.1980} or \citet{Brentari.1998}, one needs an
independent account for the agreement-data, because changes to V’s
direction or locus or handshape do not influence V’s heaviness
\citep{Brentari.1998}.

Should the aspect-construction receive an analysis independent
of the agreement-construction, then, or is there a way to unify them?
We will propose a unified analysis that takes into account both
phonological heaviness and syntactic agreement projections. We use
\citeapos{Bobaljik.Wurmbrand.2012} proposal that, for any given
derivation and its resulting LF, a combination of inviolable and
violable, but equally ranked, PF-constraints determine the best
realization of this LF. We will introduce five independently attested
PF-constraints and show that they derive the data presented here
perfectly. In this way, we demonstrate that object shift in ASL and
Libras arises from the interaction of syntax and PF, and we strengthen
current assumptions about the size of the verbal domain by making
explicit the postulated projections.

The rest of the chapter is structured as follows: In section~\ref{lasz:sec:2}, we
present our analysis in detail and show that it makes the correct
predictions. In section~\ref{lasz:sec:3}, we discuss the alternative proposals by
\citet{Fischer.Janis.1992}, \citet{Matsuoka.1997}, and \citet{Braze.2004}, who
argue for rightward-movement of V. Section~\ref{lasz:sec:4} concludes.

% ===================================================================
\section{Analysis}
    \label{lasz:sec:2}
    
% -------------------------------------------------------------------
\subsection{Phrase structure}

Our analysis extends the proposals made by \citet{Quadros.etal.2004}, 
\citet{Quadros.DLM.2010}, and \citet{Gokgoz.2013}, who follow the
common analysis of object shift in spoken languages (\citealp{Holmberg:1986}
and much work since) and suggest that in object-shift constructions,
the order of V and O is inverted by movement of O to the left. Further,
we adopt the following insights from the literature:

As stated before, we follow \citet{Fischer.1975}, \citet{Liddell.1980},
\citet{Quadros.1999} and many others in assuming that the underlying order
of ASL and Libras is SVO. In other words, VP is left-headed.

Regarding agreement with noun-classes, we adopt \citeapos{Benedicto.Brentari.2004}
proposal: For classifier-agreement with O, there is
a dedicated phrase ClP immediately dominating VP, where the
handshape-feature on V is set by V moving to the head of ClP and O
moving to Spec,ClP, followed by Agree between O and V.%
\footnote{
    The feature value that O shares with V is not the handshape feature of O’s
    sign, but O’s sign’s noun-class resp.~information about the shape of O’s referent.
}
If the
classifier expressed on V is a handling classifier, there is a second,
higher ClP responsible for agreement between V and S \citep{Benedicto.Brentari.2004}.

For agreement with loci, we adopt the proposal that loci are the
realization of morphosyntactic features, similar to the proposals by
\citet{Neidle.etal.2000}, \citet{Kuhn.2016}, and \citet{Pfau.etal.2018}. We call the
phrase where V agrees with O in locus AgrOP, and assume that it
dominates the lower ClP or, in its absence, VP.%
\footnote{
    \citet{Gokgoz.2013} argues that the phrase that objects shift to is in the
    TP-domain. He elicited word order data where in SOV sentences in ASL, manual
    sentential negation in preverbal position is rated better between O and V (SONegV)
    than before O (SNegOV): Assuming that manual negation appears in a NegP and
    that this NegP marks the edge of the VP-domain, O must raise beyond.
    
    We note that in his Experiment 3, the most accepted placement of manual
    negation is in sentence-final position for plain V, classifier-agreeing V,
    locus-agreeing V, and ditransitives agreeing in classifier and locus. (He does not
    discuss durative/continuative aspect.) We tentatively suggest that the realization of
    manual (and non-manual) negation may follow similar PF-constraints as the
    realization of O as discussed in this paper, but we consider an investigation of
    negation outside the scope of the present proposal.
    
    \citeauthor{Gokgoz.2013} also presents examples where shifted O has a presuppositional
    reading, which indicates, in line with \citet{Diesing1992}, that O moves out of the VP.
    This presuppositional effect is also observed in \citet{Napoli.SS.RMQ.2017}. We have not
    collected further data on presuppositional readings of O, but we did test for a
    definiteness effect of O in ASL (parallel to Scandinavian object-shift; \citealp{Holmberg:1986})
    and did not find one.
}
Similarly, for
directional verbs that contain movement from one locus to another, the
endpoint referring to S is established by agreement between V and S
in a projection AgrSP.

For aspect-marked verbs, we adopt parts of \citeapos{Matsuoka.1997}
and \citeapos{Braze.2004} analysis of the syntax of aspect: When the verb
displays durative-aspect agreement (reduplication), V moves to an Asp
head. However, we differ from these proposals in the headedness of
AspP: We contend that Asp is left-headed, like all verbal projections in
ASL and Libras. This allows our proposal to be consistent with the
Final-over-Final Constraint (FOFC; \citealp{BiberauerHR.2014,SheehanBRH.2017}) 
as will be discussed in section~\ref{lasz:sec:33}.

The tree in Fig.~\ref{lasz:fig:1} illustrates our assumptions about the phrase
structure of ASL and Libras for transitive verbs. Regarding
ditransitives like GIVE, which agree with the direct object in
handshape, but with the indirect object in locus, we tentatively assume
that the phrase that introduces the indirect object is projected between
the lower ClP and AgrOP. In this way, Cl will probe down and find O,
while AgrO will find IO because it is structurally higher than O.%
\footnote{
    We follow \citet{Quadros.etal.2004} in the assumption that the projections
    ClP, AgrOP, and AspP are only present when the verb shows the respective marking,
    but nothing in our analysis hinges on this.
}%
\footnote{
    See also the discussion in \citet{Pfau.etal.2018} of \citeapos{Nevins.2011} proposal.
}

\begin{figure}
    \centering
    \begin{tikzpicture}
        \Tree
        [.TP
            \node (SpecTP) {$\langle$S$\rangle$};
            [
                [.T ]
                [.Cl\textsubscript{S}P
                    \node (SpecClSP) {$\langle$S$\rangle$};
                    [
                        [.Cl\textsubscript{S} ]
                        [.AspP
                            \node (SpecAspP) {$\langle$O$\rangle$}; 
                            [
                                [.Asp ]
                                [.\textit{v}P
                                    \node (Sbase) {$\langle$S$\rangle$};
                                    [
                                        [.\textit{v} ]
                                        [.AgrOP
                                            \node (SpecAgrOP) {$\langle$O$\rangle$};
                                            [
                                                [.AgrO ]
                                                [.Cl\textsubscript{O}P
                                                    \node (SpecClOP) {$\langle$O$\rangle$};
                                                    [
                                                        [.Cl\textsubscript{O} ]
                                                        [.VP
                                                            [.V ]
                                                            \node (Obase) {$\langle$O$\rangle$};
                                                        ]
                                                    ]
                                                ]
                                            ]
                                        ]
                                    ]
                                ]
                            ]
                        ]
                    ]
                ]
            ]
        ]
        \draw[->] (Obase) [in=-90,out=-135]  to (SpecClOP);
        \draw[->] (SpecClOP) [in=-90,out=-135]  to (SpecAgrOP);
        \draw[->] (SpecAgrOP) [in=-90,out=-135]  to (SpecAspP);
        \draw[->] (Sbase) [in=-90,out=-135]  to (SpecClSP);
        \draw[->] (SpecClSP) [in=-90,out=-135]  to (SpecTP);
    \end{tikzpicture}
    \caption{Phrase structure of ASL and Libras}
    \label{lasz:fig:1}
\end{figure}

% -------------------------------------------------------------------
\subsection{Verbal agreement}
    \label{lasz:sec:22}

We make the following assumptions about verbal agreement:
In sentences where V agrees in handshape/locus/aspect, V
moves to the respective syntactic head (Cl/Agr/Asp). This movement
is covert. In other words, there is an inviolable constraint at PF stating
that it is always the lowest copy of V that will be realized. This is
similar to English, where V agrees with T, but fully lexical verbs are
never realized in T. Another example of such an inviolable constraint
is in \citeapos{Bobaljik.Wurmbrand.2012} analysis of Quantifier Raising in
English: English syntax has scrambling just like German and
Japanese, and this movement is visible at LF in all languages, but in
contrast to German and Japanese, English PF never pronounces the
higher copies, giving rise to QR.

That V does not move overtly is corroborated by the fact that
we never see V move past manual negation~\REF{lasz:ex:8} or adverbs of
frequency~\REF{lasz:ex:9} in ASL and Libras:

\ea
    \label{lasz:ex:8}
    \ea[]{ 
        IX\laszLoc{a} MARIA\laszLoc{a} NOT HELP\laszDir{a}{b} IX\laszLoc{b} JOHN\laszLoc{b}. 
        \\
        `Maria did not help John.'
        \hfill         
        \cmark Libras, 
        \cmark ASL 
        \\
        \citep[][\pnfmt{6}, ex. (26a,b)]{Quadros.etal.2004}
    }
    \ex[*]{ 
        IX\laszLoc{a} MARIA\laszLoc{a} HELP\laszDir{a}{b} NOT IX\laszLoc{b} JOHN\laszLoc{b}.
        \\
        \citep[][\pnfmt{6}, ex. (26a,b)]{Quadros.etal.2004}
        \hfill 
        *Libras, 
        *ASL 
    }
    \z 
\ex 
    \label{lasz:ex:9}
    \ea[]{
        IX\laszLoc{a} MARIA\laszLoc{a} SOMETIMES HELP\laszDir{a}{b} IX\laszLoc{b} JOHN\laszLoc{b}. 
        \hfill 
        \cmark Libras, 
        \cmark ASL
        \\
        `Maria sometimes helped John.'
        \citep[][\pnfmt{6}, ex. (27a,b)]{Quadros.etal.2004}
    }
    \ex[*]{
        IX\laszLoc{a} MARIA\laszLoc{a} HELP\laszDir{a}{b} SOMETIMES  IX\laszLoc{b} JOHN\laszLoc{b}. 
        \hfill 
        \cmark Libras, 
        \cmark ASL
        \\
        \citep[][\pnfmt{6}, ex. (27a,b)]{Quadros.etal.2004}
    }
    \z
\z 

\noindent 
When V moves to agree, this triggers syntactic movement of O to the
respective specifier position. That is, when V agrees in handshape
with O, O moves to the specifier of the low ClP (Cl\textsubscript{O}P). When V agrees in
locus with O, O moves to the specifier of AgrOP. When V agrees in
aspect, O moves to the specifier of AspP. Whether this movement of O
is overt or covert, i.e. which copy of O gets realized, depends on
PF-constraints as we will discuss below. When agreement doesn’t take
place, the agreement-phrase becomes vacuous, and we assume here
that vacuous phrases are not projected at all, but nothing hinges on this
assumption.

When V is plain, there is no position for O and V to move to.

% -------------------------------------------------------------------
\subsection{Determining PF} 
    \label{lasz:sec:23}

Syntax will thus yield a representation where V has copies in the
functional heads it agrees with, and O has copies in the specifiers of
the same functional phrases. In line with \citeapos{Bobaljik.Wurmbrand.2012} 
framework (which builds on \citealp{Bobaljik.1995,Bobaljik.2002,Brody.1995,ErteschikShir.1997}, a.o.), 
this representation is first sent to LF, and
LF decides which copies to interpret, and this will always be the
highest copies. Then possible PF-realizations of this LF are evaluated.
PF-candidates differ in which copies of V to realize (which is trivial,
given the inviolable constraint discussed above) and which copies of
O to realize. We propose the following violable and equally-ranked
constraints for PF to achieve this. We number the constraints for ease
of reference, but this numbering is not relevant for the evaluation
mechanism. Whichever PF-candidate(s) fulfil the most constraints will
be grammatical realizations for this LF.

\subsubsection*{Constraint 1: Scope Transparency}

\citet{Bobaljik.Wurmbrand.2012} investigate the connection between the
scope of quantifiers and word orders in two sets of languages: one set
allows QR and doesn’t allow scrambling (e.g. English), and the other
set allows scrambling and doesn’t allow QR (e.g. German, Japanese).
They argue that these properties are connected in the following way:
all languages allow quantifier movement in the syntax and thus at LF,
but the two sets of languages differ in which copy of the quantifier
they pronounce. Pronouncing the high copy amounts to scrambling
and is not available in English, while it is available in German and
Japanese. To capture the distribution of QR, B\&W introduce the
violable constraint Scope Transparency (10), which, given an LF with
a moved quantifier, favors pronouncing the high copy of the quantifier
(scrambling) over pronouncing the low copy of the quantifier (QR). In
English, where scrambling is not an option for PF (this is an inviolable
constraint), Scope Transparency can be violated, giving rise to QR.

\ea
    Scope Transparency (ScoT) \citep[\pnfmt{373}]{Bobaljik.Wurmbrand.2012}: \\
    If the order of two elements at LF is $A\gg B$, the order at PF is $A\gg B$. \\
    \begin{quote}
    We use the symbol $\gg$ to represent the canonical manifestation
    of hierarchical order at the relevant level: roughly, scope at LF
    [...] and linear precedence at PF [...]. \parentext{ibid.}
    \end{quote}
\z 

\citet{Bobaljik.Wurmbrand.2012} proceed to phenomena where, given
two LF choices and two PF choices, 3 out of the 4 possible
combinations are grammatical. They call this the “3⁄4 signature” and
argue that this is the characteristic pattern of the interaction of ScoT
with another violable constraint. If both constraints favor the same
PF-candidate, this PF will be the “best” PF of the given LF. If the
constraints disagree, however, this results in a tie, and both PF
candidates are the “best” choice. This pattern is illustrated in Table~\ref{lasz:tab:4}.

\begin{table}[
]
    \centering
    \begin{tabular}{ccccc}
        LF & PFs & ScoT & Constraint X & Prediction \\
        % \laszTabColA{} & 
        % \multirow{2}{*}{$A \gg B$} & 
    \end{tabular}
    \caption{The 3⁄4 signature \citep[\pnfmt{385}f.]{Bobaljik.Wurmbrand.2012}}
    \label{lasz:tab:4}
\end{table}

% ===================================================================
\printbibliography[heading=subbibliography,notkeyword=this]

\end{document}